\Page{Supports authentiques\quad /\quad DM et DS}{Travailler de vrais devoirs}{authentiques}
Les deux parcours suivants s'appuient sur des documents réellement publiés en MPSI. Ouvrir le lien pour conserver l'énoncé, ses figures et ses consignes d'origine. Les repères ci-dessous permettent de vérifier le raisonnement sans rééditer ces devoirs comme des sujets OptiTAB.
\Sub{Parcours R1\quad /\quad DM du 10 septembre 2026}
\source{MPSI, site pédagogique schleck.ovh, \href{https://mpsi.schleck.ovh/fichiers/devoirs/dm1_optique_ad.pdf}{DM1 « Optique, analyse dimensionnelle »}, exercice 2, partie II, questions 7 et 9, page 3/3.}
\textbf{Travail.} Résoudre les questions indiquées sur l'anneau catadioptrique et la rotation d'un miroir. Identifier la normale en chaque point ; utiliser le seuil verre-air avec l'indice $1{,}7$ fourni dans le devoir.
\begin{docu}{Repères de correction R1}
Le seuil vaut $\alpha_L=\arcsin(1/1{,}7)=36{,}0^\circ$. La réflexion est totale pour $\alpha>\alpha_L$, l'angle étant mesuré depuis la normale en $H$. Une rotation du miroir de $\beta$ fait tourner cette normale de $\beta$ ; la symétrie de la direction incidente par rapport à la normale fait tourner le rayon réfléchi de $2\beta$. Le point clé est la définition de l'angle, pas la mémorisation d'une figure.
\end{docu}
\Sub{Parcours R2\quad /\quad DS du 17 septembre 2026}
\source{MPSI, site pédagogique schleck.ovh, \href{https://mpsi.schleck.ovh/fichiers/devoirs/ds01_enonce.pdf}{DS1 « De l'analyse dimensionnelle et de l'optique géométrique »}, partie II « Détecteur de pluie », page 2/5. Durée du devoir complet : 2 h.}
\textbf{Travail.} Traiter la partie II avec les indices originaux du plexiglas, du verre, de l'eau et de l'air. Ne pas remplacer ce trajet d'injection par celui du DS OptiTAB.
\begin{docu}{Repères de correction R2}
À l'entrée dans le verre, $\theta_2=\arcsin[(1{,}50/1{,}55)\sin50^\circ]=47{,}85^\circ$. Les seuils verre-air et verre-eau valent $40{,}18^\circ$ et $59{,}10^\circ$. Il y a donc réflexion totale à sec, mais transmission vers l'eau ; le rayon dans l'eau fait $59{,}76^\circ$ avec la normale. L'apparition d'une voie transmise diminue la lumière réfléchie. Les lois géométriques ne chiffrent pas cette diminution.
\end{docu}
\source{Documents consultés le 25 octobre 2026. Les dates, numéros de partie et valeurs ci-dessus sont ceux des PDF sources, et non des dates créées pour ce manuel.}

\Page{Support authentique\quad /\quad Concours}{Mines-Ponts MP 2021}{auth-concours}
\source{GIP Concours commun Mines-Ponts, Physique I, filière MP, session 2021, partie IV « Observations optiques », questions 21 à 24, pages 5/6 et 6/6 du sujet.}
\begin{loi}{Accéder au sujet d'origine}
\href{https://concoursminesponts.fr/annales/}{Page officielle des annales Mines-Ponts}\par
\href{https://wordpress.concoursminesponts.fr/wp-content/uploads/2024/07/2021.zip}{Archive officielle de la session 2021}\par
Fichier dans l'archive : \texttt{2021/MP/Physique-MP-1-2021.pdf}.
\end{loi}
\Sub{Un parcours accessible après ce chapitre}
Traiter les questions 21 à 24 sur le microscope, en conservant les figures et données du sujet. Pour ce parcours ciblé, prévoir 40 minutes. La durée officielle de l'épreuve complète est de 3 heures, sans calculatrice ; cette durée ne concerne pas uniquement la partie optique.
\begin{docu}{Repères de correction R3}
\textbf{Q21.} Rayons proches de l'axe et faiblement inclinés ; stigmatisme et aplanétisme approchés.\par
\textbf{Q22.} L'image intermédiaire appartient au plan focal objet de l'oculaire ; les rayons émergents sont parallèles.\par
\textbf{Q23.} Avec $f'_1=1\ \mathrm{mm}$, $f'_2=20\ \mathrm{mm}$ et $\Delta=150\ \mathrm{mm}$, $q_1=151\ \mathrm{mm}$ et $p_1=f'_1q_1/(q_1-f'_1)=1{,}0067\ \mathrm{mm}$. L'objet est juste avant $F_1$ ; l'œil emmétrope n'accommode pas. Si l'objet est entre $F_1$ et $O_1$, son image par l'objectif devient virtuelle.\par
\textbf{Q24.} En valeur absolue, $P_i=\Delta/(f'_1f'_2)=7{,}5\times10^3\ \mathrm{m^{-1}}$. Pour le rayon de grain $R_b=0{,}2\ \mu\mathrm{m}$ donné plus haut dans le sujet, son diamètre sous-tend environ $P_i(2R_b)=3\times10^{-3}\ \mathrm{rad}$, supérieur à la résolution angulaire usuelle de l'œil. Ce modèle géométrique prédit une taille apparente suffisante.
\end{docu}
\fig{microscope}{Construction originale pour accompagner le parcours : image intermédiaire en $F_2$ et faisceau final parallèle.}
\source{Le sujet original demeure la propriété du GIP CCMP. Ses conditions de diffusion figurent sur la page officielle et dans le PDF. Ce chapitre fournit un parcours et des repères indépendants ; il ne reproduit pas l'épreuve complète.}

\Page{Références et suivi}{Un socle pour la suite du livre}{references}
\Sub{Références primaires}
\begin{enumerate}
\item Ministère de l'Enseignement supérieur, \href{https://www.education.gouv.fr/bo/21/Special1/ESRS2035779A.htm}{arrêté du 5 janvier 2021, ESRS2035779A}. \href{https://cache.media.education.gouv.fr/file/SPE1-MEN-MESRI-4-2-2021/64/8/spe779_annexe_1373648.pdf}{Annexes officielles}, annexe 2, physique-chimie MPSI ; section 1.1, pages internes 12 et 13 (pages 49 et 50 du fichier regroupé). La formation expérimentale est utilisée pour le TP et l'exploitation des mesures.
\item Site pédagogique \href{https://mpsi.schleck.ovh/}{Physique-chimie MPSI, schleck.ovh}, DM1 et DS1 de septembre 2026 ; références détaillées dans les parcours R1 et R2.
\item \href{https://concoursminesponts.fr/annales/}{GIP CCMP, annales officielles}, Physique I MP 2021, partie IV, questions 21 à 24. Les autres parties de ce sujet font appel à d'autres chapitres et ne sont pas exigées ici.
\end{enumerate}
\Sub{Nature des contenus}
Le cours, les constructions vectorielles, les 24 exercices, les trois défis, le DM et le DS OptiTAB sont rédigés pour ce chapitre. Les exemples numériques et le tableau du TP sont des données pédagogiques créées pour le travail ; ils ne sont pas présentés comme des mesures réelles. Les trois parcours R1 à R3 identifient explicitement les supports extérieurs.
\Sub{Compétences avant de poursuivre}
\checkline{Je peux refaire seul les cinq démonstrations du bilan de cours.}
\checkline{Je trace les images virtuelles aussi sûrement que les images réelles.}
\checkline{Je relie le critère de netteté à un détecteur donné.}
\checkline{Je peux corriger mon DS à partir d'une grille de capacités.}
\checkline{Je distingue un exercice d'application et un sujet authentique référencé.}
\Sub{Suite envisagée après validation}
La progression poursuivra les signaux électriques dans l'ARQS, les régimes transitoires, l'oscillateur harmonique, les régimes sinusoïdaux, le filtrage et la propagation ; puis la mécanique, la thermodynamique, l'induction et l'introduction à la physique quantique. Les outils expérimentaux et numériques seront repris à travers ces thèmes.
\begin{methode}{Version pilote}
Cette livraison permet de valider le niveau de détail, les figures, la progression et la présentation des corrigés avant de construire les chapitres suivants. Elle couvre un chapitre de physique ; elle ne constitue pas encore le programme annuel de physique-chimie.
\end{methode}
